\section{}

\paragraph{937年}\eventanchor{LIAO-0937-REGNAL-YEAR}{LIAO-0937-REGNAL-YEAR}　辽以辽太宗在位第11年作为诏令、赋役与外交文书的纪年框架。账本列出《辽史》〈本纪〉、〈营卫志〉、〈百官志〉；《宋史》辽国相关传；宋人使辽记录作为来源线索；该材料可核对时间与制度称谓，但有关执行范围、群体经验、军数或后果，仍须与地方文书、物质遗存及相关政权记录互证。

\paragraph{938年}\eventanchor{LIAO-0938-EVENT-031}{LIAO-0938-EVENT-031}　辽以幽州为南京，经营燕云城市与交通节点。账本列出《辽史》〈本纪〉、〈营卫志〉、〈百官志〉；《宋史》辽国相关传；宋人使辽记录作为来源线索；该材料可核对时间与制度称谓，但有关执行范围、群体经验、军数或后果，仍须与地方文书、物质遗存及相关政权记录互证。

\paragraph{938年}\eventanchor{LIAO-0938-REGNAL-YEAR}{LIAO-0938-REGNAL-YEAR}　辽以辽太宗在位第12年作为诏令、赋役与外交文书的纪年框架。账本列出《辽史》〈本纪〉、〈营卫志〉、〈百官志〉；《宋史》辽国相关传；宋人使辽记录作为来源线索；该材料可核对时间与制度称谓，但有关执行范围、群体经验、军数或后果，仍须与地方文书、物质遗存及相关政权记录互证。

\paragraph{939年}\eventanchor{LIAO-0939-REGNAL-YEAR}{LIAO-0939-REGNAL-YEAR}　辽以辽太宗在位第13年作为诏令、赋役与外交文书的纪年框架。账本列出《辽史》〈本纪〉、〈营卫志〉、〈百官志〉；《宋史》辽国相关传；宋人使辽记录作为来源线索；该材料可核对时间与制度称谓，但有关执行范围、群体经验、军数或后果，仍须与地方文书、物质遗存及相关政权记录互证。

\paragraph{940年}\eventanchor{LIAO-0940-REGNAL-YEAR}{LIAO-0940-REGNAL-YEAR}　辽以辽太宗在位第14年作为诏令、赋役与外交文书的纪年框架。账本列出《辽史》〈本纪〉、〈营卫志〉、〈百官志〉；《宋史》辽国相关传；宋人使辽记录作为来源线索；该材料可核对时间与制度称谓，但有关执行范围、群体经验、军数或后果，仍须与地方文书、物质遗存及相关政权记录互证。

\paragraph{941年}\eventanchor{LIAO-0941-REGNAL-YEAR}{LIAO-0941-REGNAL-YEAR}　辽以辽太宗在位第15年作为诏令、赋役与外交文书的纪年框架。账本列出《辽史》〈本纪〉、〈营卫志〉、〈百官志〉；《宋史》辽国相关传；宋人使辽记录作为来源线索；该材料可核对时间与制度称谓，但有关执行范围、群体经验、军数或后果，仍须与地方文书、物质遗存及相关政权记录互证。

\paragraph{942年}\eventanchor{LIAO-0942-REGNAL-YEAR}{LIAO-0942-REGNAL-YEAR}　辽以辽太宗在位第16年作为诏令、赋役与外交文书的纪年框架。账本列出《辽史》〈本纪〉、〈营卫志〉、〈百官志〉；《宋史》辽国相关传；宋人使辽记录作为来源线索；该材料可核对时间与制度称谓，但有关执行范围、群体经验、军数或后果，仍须与地方文书、物质遗存及相关政权记录互证。

\paragraph{943年}\eventanchor{LIAO-0943-REGNAL-YEAR}{LIAO-0943-REGNAL-YEAR}　辽以辽太宗在位第17年作为诏令、赋役与外交文书的纪年框架。账本列出《辽史》〈本纪〉、〈营卫志〉、〈百官志〉；《宋史》辽国相关传；宋人使辽记录作为来源线索；该材料可核对时间与制度称谓，但有关执行范围、群体经验、军数或后果，仍须与地方文书、物质遗存及相关政权记录互证。

\paragraph{944年}\eventanchor{LIAO-0944-REGNAL-YEAR}{LIAO-0944-REGNAL-YEAR}　辽以辽太宗在位第18年作为诏令、赋役与外交文书的纪年框架。账本列出《辽史》〈本纪〉、〈营卫志〉、〈百官志〉；《宋史》辽国相关传；宋人使辽记录作为来源线索；该材料可核对时间与制度称谓，但有关执行范围、群体经验、军数或后果，仍须与地方文书、物质遗存及相关政权记录互证。

\paragraph{945年}\eventanchor{LIAO-0945-REGNAL-YEAR}{LIAO-0945-REGNAL-YEAR}　辽以辽太宗在位第19年作为诏令、赋役与外交文书的纪年框架。账本列出《辽史》〈本纪〉、〈营卫志〉、〈百官志〉；《宋史》辽国相关传；宋人使辽记录作为来源线索；该材料可核对时间与制度称谓，但有关执行范围、群体经验、军数或后果，仍须与地方文书、物质遗存及相关政权记录互证。

\paragraph{946年}\eventanchor{LIAO-0946-REGNAL-YEAR}{LIAO-0946-REGNAL-YEAR}　辽以辽太宗在位第20年作为诏令、赋役与外交文书的纪年框架。账本列出《辽史》〈本纪〉、〈营卫志〉、〈百官志〉；《宋史》辽国相关传；宋人使辽记录作为来源线索；该材料可核对时间与制度称谓，但有关执行范围、群体经验、军数或后果，仍须与地方文书、物质遗存及相关政权记录互证。

\paragraph{947年}\eventanchor{LIAO-0947-EVENT-032}{LIAO-0947-EVENT-032}　辽军进入开封，后因补给与抵抗压力撤离。账本列出《辽史》〈本纪〉、〈营卫志〉、〈百官志〉；《宋史》辽国相关传；宋人使辽记录作为来源线索；该材料可核对时间与制度称谓，但有关执行范围、群体经验、军数或后果，仍须与地方文书、物质遗存及相关政权记录互证。

\paragraph{947年}\eventanchor{LIAO-0947-REGNAL-YEAR}{LIAO-0947-REGNAL-YEAR}　辽以辽太宗在位第21年作为诏令、赋役与外交文书的纪年框架。账本列出《辽史》〈本纪〉、〈营卫志〉、〈百官志〉；《宋史》辽国相关传；宋人使辽记录作为来源线索；该材料可核对时间与制度称谓，但有关执行范围、群体经验、军数或后果，仍须与地方文书、物质遗存及相关政权记录互证。

\paragraph{948年}\eventanchor{LIAO-0948-REGNAL-YEAR}{LIAO-0948-REGNAL-YEAR}　辽以辽世宗在位第1年作为诏令、赋役与外交文书的纪年框架。账本列出《辽史》〈本纪〉、〈营卫志〉、〈百官志〉；《宋史》辽国相关传；宋人使辽记录作为来源线索；该材料可核对时间与制度称谓，但有关执行范围、群体经验、军数或后果，仍须与地方文书、物质遗存及相关政权记录互证。

\paragraph{949年}\eventanchor{LIAO-0949-REGNAL-YEAR}{LIAO-0949-REGNAL-YEAR}　辽以辽世宗在位第2年作为诏令、赋役与外交文书的纪年框架。账本列出《辽史》〈本纪〉、〈营卫志〉、〈百官志〉；《宋史》辽国相关传；宋人使辽记录作为来源线索；该材料可核对时间与制度称谓，但有关执行范围、群体经验、军数或后果，仍须与地方文书、物质遗存及相关政权记录互证。

\paragraph{950年}\eventanchor{LIAO-0950-REGNAL-YEAR}{LIAO-0950-REGNAL-YEAR}　辽以辽世宗在位第3年作为诏令、赋役与外交文书的纪年框架。账本列出《辽史》〈本纪〉、〈营卫志〉、〈百官志〉；《宋史》辽国相关传；宋人使辽记录作为来源线索；该材料可核对时间与制度称谓，但有关执行范围、群体经验、军数或后果，仍须与地方文书、物质遗存及相关政权记录互证。

\paragraph{951年}\eventanchor{LIAO-0951-REGNAL-YEAR}{LIAO-0951-REGNAL-YEAR}　辽以辽世宗在位第4年作为诏令、赋役与外交文书的纪年框架。账本列出《辽史》〈本纪〉、〈营卫志〉、〈百官志〉；《宋史》辽国相关传；宋人使辽记录作为来源线索；该材料可核对时间与制度称谓，但有关执行范围、群体经验、军数或后果，仍须与地方文书、物质遗存及相关政权记录互证。

\paragraph{952年}\eventanchor{LIAO-0952-REGNAL-YEAR}{LIAO-0952-REGNAL-YEAR}　辽以辽穆宗在位第1年作为诏令、赋役与外交文书的纪年框架。账本列出《辽史》〈本纪〉、〈营卫志〉、〈百官志〉；《宋史》辽国相关传；宋人使辽记录作为来源线索；该材料可核对时间与制度称谓，但有关执行范围、群体经验、军数或后果，仍须与地方文书、物质遗存及相关政权记录互证。

\paragraph{953年}\eventanchor{LIAO-0953-REGNAL-YEAR}{LIAO-0953-REGNAL-YEAR}　辽以辽穆宗在位第2年作为诏令、赋役与外交文书的纪年框架。账本列出《辽史》〈本纪〉、〈营卫志〉、〈百官志〉；《宋史》辽国相关传；宋人使辽记录作为来源线索；该材料可核对时间与制度称谓，但有关执行范围、群体经验、军数或后果，仍须与地方文书、物质遗存及相关政权记录互证。

\paragraph{954年}\eventanchor{LIAO-0954-REGNAL-YEAR}{LIAO-0954-REGNAL-YEAR}　辽以辽穆宗在位第3年作为诏令、赋役与外交文书的纪年框架。账本列出《辽史》〈本纪〉、〈营卫志〉、〈百官志〉；《宋史》辽国相关传；宋人使辽记录作为来源线索；该材料可核对时间与制度称谓，但有关执行范围、群体经验、军数或后果，仍须与地方文书、物质遗存及相关政权记录互证。

\paragraph{955年}\eventanchor{LIAO-0955-REGNAL-YEAR}{LIAO-0955-REGNAL-YEAR}　辽以辽穆宗在位第4年作为诏令、赋役与外交文书的纪年框架。账本列出《辽史》〈本纪〉、〈营卫志〉、〈百官志〉；《宋史》辽国相关传；宋人使辽记录作为来源线索；该材料可核对时间与制度称谓，但有关执行范围、群体经验、军数或后果，仍须与地方文书、物质遗存及相关政权记录互证。

\paragraph{956年}\eventanchor{LIAO-0956-REGNAL-YEAR}{LIAO-0956-REGNAL-YEAR}　辽以辽穆宗在位第5年作为诏令、赋役与外交文书的纪年框架。账本列出《辽史》〈本纪〉、〈营卫志〉、〈百官志〉；《宋史》辽国相关传；宋人使辽记录作为来源线索；该材料可核对时间与制度称谓，但有关执行范围、群体经验、军数或后果，仍须与地方文书、物质遗存及相关政权记录互证。

\paragraph{957年}\eventanchor{LIAO-0957-REGNAL-YEAR}{LIAO-0957-REGNAL-YEAR}　辽以辽穆宗在位第6年作为诏令、赋役与外交文书的纪年框架。账本列出《辽史》〈本纪〉、〈营卫志〉、〈百官志〉；《宋史》辽国相关传；宋人使辽记录作为来源线索；该材料可核对时间与制度称谓，但有关执行范围、群体经验、军数或后果，仍须与地方文书、物质遗存及相关政权记录互证。

\paragraph{958年}\eventanchor{LIAO-0958-REGNAL-YEAR}{LIAO-0958-REGNAL-YEAR}　辽以辽穆宗在位第7年作为诏令、赋役与外交文书的纪年框架。账本列出《辽史》〈本纪〉、〈营卫志〉、〈百官志〉；《宋史》辽国相关传；宋人使辽记录作为来源线索；该材料可核对时间与制度称谓，但有关执行范围、群体经验、军数或后果，仍须与地方文书、物质遗存及相关政权记录互证。

\paragraph{959年}\eventanchor{LIAO-0959-REGNAL-YEAR}{LIAO-0959-REGNAL-YEAR}　辽以辽穆宗在位第8年作为诏令、赋役与外交文书的纪年框架。账本列出《辽史》〈本纪〉、〈营卫志〉、〈百官志〉；《宋史》辽国相关传；宋人使辽记录作为来源线索；该材料可核对时间与制度称谓，但有关执行范围、群体经验、军数或后果，仍须与地方文书、物质遗存及相关政权记录互证。

\paragraph{960年}\eventanchor{LIAO-0960-REGNAL-YEAR}{LIAO-0960-REGNAL-YEAR}　辽以辽穆宗在位第9年作为诏令、赋役与外交文书的纪年框架。账本列出《辽史》〈本纪〉、〈营卫志〉、〈百官志〉；《宋史》辽国相关传；宋人使辽记录作为来源线索；该材料可核对时间与制度称谓，但有关执行范围、群体经验、军数或后果，仍须与地方文书、物质遗存及相关政权记录互证。

\paragraph{961年}\eventanchor{LIAO-0961-REGNAL-YEAR}{LIAO-0961-REGNAL-YEAR}　辽以辽穆宗在位第10年作为诏令、赋役与外交文书的纪年框架。账本列出《辽史》〈本纪〉、〈营卫志〉、〈百官志〉；《宋史》辽国相关传；宋人使辽记录作为来源线索；该材料可核对时间与制度称谓，但有关执行范围、群体经验、军数或后果，仍须与地方文书、物质遗存及相关政权记录互证。

\paragraph{962年}\eventanchor{LIAO-0962-REGNAL-YEAR}{LIAO-0962-REGNAL-YEAR}　辽以辽穆宗在位第11年作为诏令、赋役与外交文书的纪年框架。账本列出《辽史》〈本纪〉、〈营卫志〉、〈百官志〉；《宋史》辽国相关传；宋人使辽记录作为来源线索；该材料可核对时间与制度称谓，但有关执行范围、群体经验、军数或后果，仍须与地方文书、物质遗存及相关政权记录互证。

\paragraph{963年}\eventanchor{LIAO-0963-REGNAL-YEAR}{LIAO-0963-REGNAL-YEAR}　辽以辽穆宗在位第12年作为诏令、赋役与外交文书的纪年框架。账本列出《辽史》〈本纪〉、〈营卫志〉、〈百官志〉；《宋史》辽国相关传；宋人使辽记录作为来源线索；该材料可核对时间与制度称谓，但有关执行范围、群体经验、军数或后果，仍须与地方文书、物质遗存及相关政权记录互证。

\paragraph{964年}\eventanchor{LIAO-0964-REGNAL-YEAR}{LIAO-0964-REGNAL-YEAR}　辽以辽穆宗在位第13年作为诏令、赋役与外交文书的纪年框架。账本列出《辽史》〈本纪〉、〈营卫志〉、〈百官志〉；《宋史》辽国相关传；宋人使辽记录作为来源线索；该材料可核对时间与制度称谓，但有关执行范围、群体经验、军数或后果，仍须与地方文书、物质遗存及相关政权记录互证。

\paragraph{965年}\eventanchor{LIAO-0965-REGNAL-YEAR}{LIAO-0965-REGNAL-YEAR}　辽以辽穆宗在位第14年作为诏令、赋役与外交文书的纪年框架。账本列出《辽史》〈本纪〉、〈营卫志〉、〈百官志〉；《宋史》辽国相关传；宋人使辽记录作为来源线索；该材料可核对时间与制度称谓，但有关执行范围、群体经验、军数或后果，仍须与地方文书、物质遗存及相关政权记录互证。

\paragraph{966年}\eventanchor{LIAO-0966-REGNAL-YEAR}{LIAO-0966-REGNAL-YEAR}　辽以辽穆宗在位第15年作为诏令、赋役与外交文书的纪年框架。账本列出《辽史》〈本纪〉、〈营卫志〉、〈百官志〉；《宋史》辽国相关传；宋人使辽记录作为来源线索；该材料可核对时间与制度称谓，但有关执行范围、群体经验、军数或后果，仍须与地方文书、物质遗存及相关政权记录互证。

\paragraph{967年}\eventanchor{LIAO-0967-REGNAL-YEAR}{LIAO-0967-REGNAL-YEAR}　辽以辽穆宗在位第16年作为诏令、赋役与外交文书的纪年框架。账本列出《辽史》〈本纪〉、〈营卫志〉、〈百官志〉；《宋史》辽国相关传；宋人使辽记录作为来源线索；该材料可核对时间与制度称谓，但有关执行范围、群体经验、军数或后果，仍须与地方文书、物质遗存及相关政权记录互证。

\paragraph{968年}\eventanchor{LIAO-0968-REGNAL-YEAR}{LIAO-0968-REGNAL-YEAR}　辽以辽穆宗在位第17年作为诏令、赋役与外交文书的纪年框架。账本列出《辽史》〈本纪〉、〈营卫志〉、〈百官志〉；《宋史》辽国相关传；宋人使辽记录作为来源线索；该材料可核对时间与制度称谓，但有关执行范围、群体经验、军数或后果，仍须与地方文书、物质遗存及相关政权记录互证。

\paragraph{969年}\eventanchor{LIAO-0969-EVENT-033}{LIAO-0969-EVENT-033}　辽景宗即位，萧绰逐步取得政治作用。账本列出《辽史》〈本纪〉、〈营卫志〉、〈百官志〉；《宋史》辽国相关传；宋人使辽记录作为来源线索；该材料可核对时间与制度称谓，但有关执行范围、群体经验、军数或后果，仍须与地方文书、物质遗存及相关政权记录互证。

\paragraph{969年}\eventanchor{LIAO-0969-REGNAL-YEAR}{LIAO-0969-REGNAL-YEAR}　辽以辽穆宗在位第18年作为诏令、赋役与外交文书的纪年框架。账本列出《辽史》〈本纪〉、〈营卫志〉、〈百官志〉；《宋史》辽国相关传；宋人使辽记录作为来源线索；该材料可核对时间与制度称谓，但有关执行范围、群体经验、军数或后果，仍须与地方文书、物质遗存及相关政权记录互证。

\paragraph{970年}\eventanchor{LIAO-0970-REGNAL-YEAR}{LIAO-0970-REGNAL-YEAR}　辽以辽景宗在位第1年作为诏令、赋役与外交文书的纪年框架。账本列出《辽史》〈本纪〉、〈营卫志〉、〈百官志〉；《宋史》辽国相关传；宋人使辽记录作为来源线索；该材料可核对时间与制度称谓，但有关执行范围、群体经验、军数或后果，仍须与地方文书、物质遗存及相关政权记录互证。

\paragraph{971年}\eventanchor{LIAO-0971-REGNAL-YEAR}{LIAO-0971-REGNAL-YEAR}　辽以辽景宗在位第2年作为诏令、赋役与外交文书的纪年框架。账本列出《辽史》〈本纪〉、〈营卫志〉、〈百官志〉；《宋史》辽国相关传；宋人使辽记录作为来源线索；该材料可核对时间与制度称谓，但有关执行范围、群体经验、军数或后果，仍须与地方文书、物质遗存及相关政权记录互证。

\paragraph{972年}\eventanchor{LIAO-0972-REGNAL-YEAR}{LIAO-0972-REGNAL-YEAR}　辽以辽景宗在位第3年作为诏令、赋役与外交文书的纪年框架。账本列出《辽史》〈本纪〉、〈营卫志〉、〈百官志〉；《宋史》辽国相关传；宋人使辽记录作为来源线索；该材料可核对时间与制度称谓，但有关执行范围、群体经验、军数或后果，仍须与地方文书、物质遗存及相关政权记录互证。

\paragraph{973年}\eventanchor{LIAO-0973-REGNAL-YEAR}{LIAO-0973-REGNAL-YEAR}　辽以辽景宗在位第4年作为诏令、赋役与外交文书的纪年框架。账本列出《辽史》〈本纪〉、〈营卫志〉、〈百官志〉；《宋史》辽国相关传；宋人使辽记录作为来源线索；该材料可核对时间与制度称谓，但有关执行范围、群体经验、军数或后果，仍须与地方文书、物质遗存及相关政权记录互证。

\paragraph{974年}\eventanchor{LIAO-0974-REGNAL-YEAR}{LIAO-0974-REGNAL-YEAR}　辽以辽景宗在位第5年作为诏令、赋役与外交文书的纪年框架。账本列出《辽史》〈本纪〉、〈营卫志〉、〈百官志〉；《宋史》辽国相关传；宋人使辽记录作为来源线索；该材料可核对时间与制度称谓，但有关执行范围、群体经验、军数或后果，仍须与地方文书、物质遗存及相关政权记录互证。

\paragraph{975年}\eventanchor{LIAO-0975-REGNAL-YEAR}{LIAO-0975-REGNAL-YEAR}　辽以辽景宗在位第6年作为诏令、赋役与外交文书的纪年框架。账本列出《辽史》〈本纪〉、〈营卫志〉、〈百官志〉；《宋史》辽国相关传；宋人使辽记录作为来源线索；该材料可核对时间与制度称谓，但有关执行范围、群体经验、军数或后果，仍须与地方文书、物质遗存及相关政权记录互证。

\paragraph{976年}\eventanchor{LIAO-0976-REGNAL-YEAR}{LIAO-0976-REGNAL-YEAR}　辽以辽景宗在位第7年作为诏令、赋役与外交文书的纪年框架。账本列出《辽史》〈本纪〉、〈营卫志〉、〈百官志〉；《宋史》辽国相关传；宋人使辽记录作为来源线索；该材料可核对时间与制度称谓，但有关执行范围、群体经验、军数或后果，仍须与地方文书、物质遗存及相关政权记录互证。

\paragraph{977年}\eventanchor{LIAO-0977-REGNAL-YEAR}{LIAO-0977-REGNAL-YEAR}　辽以辽景宗在位第8年作为诏令、赋役与外交文书的纪年框架。账本列出《辽史》〈本纪〉、〈营卫志〉、〈百官志〉；《宋史》辽国相关传；宋人使辽记录作为来源线索；该材料可核对时间与制度称谓，但有关执行范围、群体经验、军数或后果，仍须与地方文书、物质遗存及相关政权记录互证。

\paragraph{978年}\eventanchor{LIAO-0978-REGNAL-YEAR}{LIAO-0978-REGNAL-YEAR}　辽以辽景宗在位第9年作为诏令、赋役与外交文书的纪年框架。账本列出《辽史》〈本纪〉、〈营卫志〉、〈百官志〉；《宋史》辽国相关传；宋人使辽记录作为来源线索；该材料可核对时间与制度称谓，但有关执行范围、群体经验、军数或后果，仍须与地方文书、物质遗存及相关政权记录互证。

\paragraph{979年}\eventanchor{LIAO-0979-EVENT-034}{LIAO-0979-EVENT-034}　辽军在高梁河击退北伐宋军。账本列出《辽史》〈本纪〉、〈营卫志〉、〈百官志〉；《宋史》辽国相关传；宋人使辽记录作为来源线索；该材料可核对时间与制度称谓，但有关执行范围、群体经验、军数或后果，仍须与地方文书、物质遗存及相关政权记录互证。

\paragraph{979年}\eventanchor{LIAO-0979-REGNAL-YEAR}{LIAO-0979-REGNAL-YEAR}　辽以辽景宗在位第10年作为诏令、赋役与外交文书的纪年框架。账本列出《辽史》〈本纪〉、〈营卫志〉、〈百官志〉；《宋史》辽国相关传；宋人使辽记录作为来源线索；该材料可核对时间与制度称谓，但有关执行范围、群体经验、军数或后果，仍须与地方文书、物质遗存及相关政权记录互证。

\paragraph{980年}\eventanchor{LIAO-0980-REGNAL-YEAR}{LIAO-0980-REGNAL-YEAR}　辽以辽景宗在位第11年作为诏令、赋役与外交文书的纪年框架。账本列出《辽史》〈本纪〉、〈营卫志〉、〈百官志〉；《宋史》辽国相关传；宋人使辽记录作为来源线索；该材料可核对时间与制度称谓，但有关执行范围、群体经验、军数或后果，仍须与地方文书、物质遗存及相关政权记录互证。

\paragraph{981年}\eventanchor{LIAO-0981-REGNAL-YEAR}{LIAO-0981-REGNAL-YEAR}　辽以辽景宗在位第12年作为诏令、赋役与外交文书的纪年框架。账本列出《辽史》〈本纪〉、〈营卫志〉、〈百官志〉；《宋史》辽国相关传；宋人使辽记录作为来源线索；该材料可核对时间与制度称谓，但有关执行范围、群体经验、军数或后果，仍须与地方文书、物质遗存及相关政权记录互证。

\paragraph{982年}\eventanchor{LIAO-0982-EVENT-035}{LIAO-0982-EVENT-035}　辽圣宗即位，萧太后摄政。账本列出《辽史》〈本纪〉、〈营卫志〉、〈百官志〉；《宋史》辽国相关传；宋人使辽记录作为来源线索；该材料可核对时间与制度称谓，但有关执行范围、群体经验、军数或后果，仍须与地方文书、物质遗存及相关政权记录互证。

